%! Rates of Change in Polar Functions — Unit 3, Lesson 3.15 — Exit Ticket (Answer Key)
\documentclass[10pt]{article}
\usepackage{precalculus-article}
\usepackage{precalculus-key}   % pulls in -boxes; do NOT also load -boxes
\newcommand{\TallMath}[1]{$\displaystyle #1\rule[-1.4em]{0pt}{3.2em}$}
\newcommand{\CourseName}{Precalculus}
\newcommand{\SchoolYear}{2026--2027}
\newcommand{\MeetingLength}{55 minutes}
\newcommand{\UnitNumberName}{Unit 3: Trigonometric and Polar Functions \quad}
\newcommand{\LessonNumberName}{Lesson 3.15: Rates of Change in Polar Functions}

\begin{document}

\pageheader{Unit 3, Lesson 3.15}{Exit Ticket --- Answer Key}
\namedateperiod

\vspace{0.12in}

\begin{notesbox}{Exit Ticket --- Key}

\textbf{1.}\quad The table below gives partial values of $r = 1 + \sin\theta$.

\vspace{6pt}
\renewcommand{\arraystretch}{1.5}
\begin{tabular}{|c|c|c|c|c|}
\hline
$\theta$ & $0$ & $\pi/2$ & $\pi$ & $3\pi/2$ \\
\hline
$r = 1 + \sin\theta$ & $1$ & $2$ & $1$ & $0$ \\
\hline
\end{tabular}

\vspace{8pt}
(a)\quad On $[0, \pi/2]$: \ans{Positive and increasing} ($r$ goes from 1 to 2).

\vspace{8pt}
(b)\quad On $[\pi/2, \pi]$: \ans{Positive and decreasing} ($r$ goes from 2 to 1).

\vspace{0.16in}
\textbf{2.}\quad Average rate of change from $\theta = \pi/2$ to $\theta = \pi$:

$$\frac{r(\pi) - r(\pi/2)}{\pi - \pi/2} = \frac{{\color{keyred}\mathbf{1}} - {\color{keyred}\mathbf{2}}}{\pi/2} = \ans{$-2/\pi \approx -0.64$ per radian}$$

\vspace{0.16in}
\textbf{3.}\quad \ans{A negative average rate of change of $r$ means that $r$ is decreasing over the interval --- as $\theta$ increases, the radius (distance from the origin) is getting smaller, so the curve is tracing toward the origin (if $r > 0$).}

\end{notesbox}

\end{document}
